% Irodov, user-provided Arihant PDF, printed pp.56-57, PDF pp.65-66.
% Relativity Section 1.8. The original answer for #349 prints 2.2 x 10^4
% m/s, contradicting both its formula and the data; derived value ~2.2 x 10^8.
\exproblem{340}{56}
一根细杆沿自身长度方向以恒速 $v$
相对惯性参考系 $K$ 运动。
问速度为多少时，$K$ 系中测得的杆长
比其固有长度减少 $\eta=0.5\%$？
\exsolution
记杆固有长度 $L_0$，
运动方向的洛伦兹收缩为
\[
L=L_0\sqrt{1-\beta^2},\qquad\beta=v/c.
\]
长度比固有长度减少比例
$\eta=(L_0-L)/L_0$，故
$\sqrt{1-\beta^2}=1-\eta$。
两边平方：
\[
\beta^2=1-(1-\eta)^2=\eta(2-\eta).
\]
因此
\[
\boxed{v=c\sqrt{\eta(2-\eta)}
=c\sqrt{0.005(1.995)}
\simeq0.100c}.
\]
\excheck 相对速度虽然约为光速的 $10\%$，
长度收缩仍只有 $0.5\%$，因为低速收缩是
$\beta^2/2$ 量级。

\exproblem{341}{56}
一个等边三角形三边的固有长度都是 $a$。
在相对该三角形静止系以速率 $V$
作匀速运动的惯性系中，
若相对运动方向分别沿三角形
(a) 一条角平分线；(b) 一条边，
求三角形\emph{同一时刻测得}的周长；
并考察 $V\ll c$ 和 $V\to c$ 两种极限。
\exsolution
令 $\beta=V/c$，
沿相对运动方向的长度分量收缩因子
$q=\sqrt{1-\beta^2}$，垂直分量保持不变。
一条静止时与运动方向夹角为 $\theta$ 的边，
运动系长度为
\[
a'(\theta)=a\sqrt{q^2\cos^2\theta+\sin^2\theta}.
\]
(a) 角平分线方向与底边垂直，底边仍长 $a$；
两条腰与此方向分别成 $30^\circ$，
各腰测长为
$a\sqrt{(3/4)q^2+1/4}$。
于是
\[
\boxed{P_a=a\left[1+\sqrt{4-3\beta^2}\right]}.
\]
(b) 一条边平行于相对运动方向，
测长 $aq$；另两边各与运动方向成 $60^\circ$，
各测长 $a\sqrt{q^2/4+3/4}$，所以
\[
\boxed{P_b=a\left[
\sqrt{1-\beta^2}+\sqrt{4-\beta^2}\right]}.
\]
低速展开二者均为
$P=3a-\tfrac34a\beta^2+O(\beta^4)$。
当 $\beta\to1^-$ 时，
$\boxed{P_a\to2a,\quad P_b\to\sqrt3a}$。
\excheck 对每条边测量须在\emph{运动系同一时刻}
取两端坐标；不能直接把固有周长整体乘同一个
洛伦兹因子。

\exproblem{342}{56}
一根细杆在实验室系中以速度
$v=c/2$ 平动，其测得长度为
$L=1.00\,\mathrm m$；
杆与其速度方向的夹角
$\theta=45^\circ$。
求杆的固有长度 $L_0$。
\exsolution
在实验室中把杆长分解为
$L_\parallel=L\cos\theta$ 与
$L_\perp=L\sin\theta$。
洛伦兹收缩只发生在平行运动方向，
因此在杆静止系中
$L_{0,\parallel}=\gamma L_\parallel$，
$L_{0,\perp}=L_\perp$，
其中 $\gamma=1/\sqrt{1-\beta^2}$。
所以
\[
L_0^2=\gamma^2L^2\cos^2\theta+L^2\sin^2\theta
=L^2\frac{1-\beta^2\sin^2\theta}{1-\beta^2}.
\]
故
\[
\boxed{L_0=L\sqrt{
\frac{1-\beta^2\sin^2\theta}{1-\beta^2}}
=\sqrt{\frac{1-\frac14\cdot\frac12}{1-\frac14}}
\,\mathrm m\simeq1.08\,\mathrm m}.
\]
\excheck 原书答案排印中平方根结构可能不易辨认；
只有这一表达式满足长度量纲与平、垂直分量关系。

\exproblem{343}{56}
静止的直立圆锥顶角参数为
$\theta=45^\circ$（母线与轴线夹角），
侧面积 $S_0=4.0\,\mathrm{m^2}$。
在相对它以速度 $v=(4/5)c$
沿锥轴平移的惯性系中，求：
(a) 母线与轴线夹角 $\theta'$；
(b) 锥体侧面积 $S'$。
\exsolution
设锥底半径为 $R$、固有高度为 $H$，
于是 $\tan\theta=R/H$，
固有母线长 $s_0=\sqrt{R^2+H^2}$，
侧面积 $S_0=\pi R s_0$。
运动系中轴向高度收缩为
$H'=H\sqrt{1-\beta^2}$，
横向半径仍为 $R$。
因此
\[
\tan\theta'=\frac R{H'}
=\frac{\tan\theta}{\sqrt{1-\beta^2}},
\]
$\beta=4/5$、$\theta=45^\circ$ 时
$\tan\theta'=5/3$，故
$\boxed{\theta'\simeq59^\circ}$。
新的侧面积
$S'=\pi R\sqrt{R^2+H'^2}$，
以 $S_0$ 表示：
\[
\boxed{S'=S_0
\sqrt{1-\beta^2\cos^2\theta}
=4.0\sqrt{1-0.64(1/2)}
\simeq3.3\,\mathrm{m^2}}.
\]
\excheck 曲面面积不简单地统一乘以
$\sqrt{1-\beta^2}$；半径方向不收缩，
所以须重算母线长。

\exproblem{344}{56}
一只钟相对惯性系 $K$ 匀速运动。
$K$ 系中的标准钟经过
$t=5.0\,\mathrm s$ 时，
移动的钟只走过 $t-\Delta t$，
其中慢了 $\Delta t=0.10\,\mathrm s$。
求移动钟的速率 $v$。
\exsolution
移动钟所记录的是其世界线上的固有时间
$\tau=t-\Delta t$，时间膨胀关系为
\[
\tau=t\sqrt{1-v^2/c^2}.
\]
故
\[
\sqrt{1-\beta^2}=1-\frac{\Delta t}{t},
\quad
\beta=\sqrt{1-\left(1-\frac{\Delta t}{t}\right)^2}.
\]
代入 $\Delta t/t=0.10/5.0=0.020$，
\[
\boxed{v=c\sqrt{\frac{\Delta t}{t}
\left(2-\frac{\Delta t}{t}\right)}
\simeq0.199c
\simeq6.0\times10^7\,\mathrm{m/s}}.
\]
\excheck 必须比较同一对事件的运动钟固有时间和
$K$ 系坐标时间，不是“两只静止钟有不同走时率”。

\exproblem{345}{56}
一根杆沿自身长度方向相对于地面惯性系 $K$
以恒速飞过地面一个固定标记。
地面系中从杆头至杆尾通过标记历时
$\Delta t=20\,\mathrm{ns}$；
杆的静止系中，该标记通过整根杆历时
$\Delta t'=25\,\mathrm{ns}$。
求杆的固有长度 $L_0$。
\exsolution
记杆的速率为 $v$、洛伦兹因子 $\gamma$。
在 $K$ 中杆的长度收缩为 $L=L_0/\gamma$，
经过固定点所需时间
$\Delta t=L/v=L_0/(\gamma v)$。
在杆静止系中标记以速率 $v$
相对杆运动，走过固有长度所需
$\Delta t'=L_0/v$。
因此 $\gamma=\Delta t'/\Delta t=25/20=1.25$，
\[
\beta=\sqrt{1-\gamma^{-2}}
=\sqrt{1-\frac1{1.25^2}}=0.60.
\]
杆的固有长度
\[
\boxed{L_0=v\Delta t'
=(0.60c)(25\times10^{-9}\,\mathrm s)
\simeq4.5\,\mathrm m}.
\]
\excheck 两个时间间隔是在不同参考系测量的
“整根杆通过一个静止标记”的时间，
不可直接认为二者相等。

\exproblem{346}{56}
一颗不稳定粒子的固有寿命为
$\Delta\tau_0=10\,\mathrm{ns}$，
实验室系中从产生到衰变的时间间隔
$\Delta t=20\,\mathrm{ns}$。
求实验室中粒子在寿命内飞过的距离。
\exsolution
产生与衰变是粒子世界线上的两个事件，
粒子本身固有时间为 $\Delta\tau_0$。
相对实验室的洛伦兹因子
$\gamma=\Delta t/\Delta\tau_0=2$，
速度大小
\[
v=c\sqrt{1-\gamma^{-2}}
=c\frac{\sqrt3}{2}.
\]
实验室飞行距离为
\[
\boxed{L=v\Delta t
=c\Delta t\sqrt{
1-(\Delta\tau_0/\Delta t)^2}
\simeq5.2\,\mathrm m\approx5\,\mathrm m}.
\]
\excheck 固有寿命不是实验室中读出的寿命，
必须使用时间膨胀关系求粒子速度。

\exproblem{347}{56}
一颗 $\mu$ 子在实验室系中以
$v=0.990c$ 运动，
从产生点飞到衰变点的实验室距离为
$\ell=3.0\,\mathrm{km}$。
求：(a) $\mu$ 子的固有寿命；
(b) 在其自身静止系中，产生点与衰变点的相关
参照距离（等效为实验室中两固定位置的长度收缩值）。
\exsolution
在实验室中两事件相隔的坐标时间为
$\Delta t=\ell/v$。
粒子沿自身世界线经过的固有时间为
\[
\boxed{\Delta\tau_0=
\Delta t\sqrt{1-\beta^2}
=\frac{\ell}{v}\sqrt{1-\beta^2}}.
\]
代入 $\beta=0.990$，
$\sqrt{1-\beta^2}=\sqrt{0.0199}\simeq0.1411$，
得到 $\boxed{\Delta\tau_0\simeq1.43\,\mathrm{\mu s}}$。
在 $\mu$ 子系中，实验室的产生、衰变位置间
固有实验室尺距 $\ell$ 被收缩为
\[
\boxed{\ell'=\ell\sqrt{1-\beta^2}
\simeq0.423\,\mathrm{km}\approx0.42\,\mathrm{km}}.
\]
\excheck 在粒子自身系，粒子在产生与衰变之间
没有空间位移；上述 $\ell'$ 是\emph{移动的实验室两位置}
在粒子系中同一时刻测得的距离。
必须区分这一尺距与事件位移的概念。

\exproblem{348}{56}
两颗粒子在实验室中沿同一直线同向以
$v=(3/4)c$ 匀速运动，
先后击中同一个固定靶，两次击中
在实验室中相隔
$\Delta t=50\,\mathrm{ns}$。
求两粒子在其共同静止系中的固有间距。
\exsolution
两颗粒子在实验室中保持恒定间距
$L=v\Delta t$，因为同一靶被两者先后击中，
靶在两次碰撞间静止、相对粒子移动距离为 $v\Delta t$。
粒子共同静止系中的间距为固有长度，
实验室中沿运动方向收缩：
$L=L_0/\gamma$。
所以
\[
\boxed{L_0=\gamma v\Delta t
=\frac{v\Delta t}{\sqrt{1-v^2/c^2}}}.
\]
代入 $\beta=3/4$，
$\gamma=4/\sqrt7$，得到
$\boxed{L_0\simeq17\,\mathrm m}$。
\excheck 固有间距应比实验室中的
$v\Delta t\simeq11.25\,\mathrm m$ 更大。

\exproblem{349}{56}
一根杆沿一把静止尺匀速运动。
若杆两端坐标在\emph{尺的静止系同时}读取，
两刻度差为 $\Delta x_1=4.0\,\mathrm m$；
若杆两端坐标在\emph{杆自身静止系同时}读取，
读到的仍是尺刻度，但两位置对应的刻度差为
$\Delta x_2=9.0\,\mathrm m$。
求杆固有长度 $L_0$ 与相对尺的速率 $v$。
\exsolution
在尺静止系的同时测量给出长度收缩
$\Delta x_1=L_0/\gamma$。
在杆静止系的同时测量要求
$\Delta t'=0$，由逆洛伦兹变换，
这两个测量事件在尺系满足
$\Delta x=\gamma(\Delta x'+v\Delta t')
=\gamma L_0$，故
$\Delta x_2=\gamma L_0$。
联立得到
\[
\boxed{L_0=\sqrt{\Delta x_1\Delta x_2}
=\sqrt{4.0\times9.0}=6.0\,\mathrm m},
\]
\[
\gamma=\sqrt{\frac{\Delta x_2}{\Delta x_1}}
=\frac32,\qquad
\boxed{v=c\sqrt{1-\frac{\Delta x_1}{\Delta x_2}}
=c\sqrt{\frac59}\simeq2.24\times10^8\,\mathrm{m/s}}.
\]
\excheck \textbf{原书数值勘误：}本版答案页同时印有上述
正确公式和数值“$2.2\times10^4\,\mathrm{m/s}$”，
但 $c\sqrt{5/9}\simeq2.2\times10^8\,\mathrm{m/s}$，
相差 $10^4$ 倍；按公式和题面数据的结果已明确给出。

\exproblem{350}{57}
两根固有长度均为 $L_0$ 的杆沿共同轴线
彼此相向匀速运动。
在其中一根杆的静止系中，
两杆\emph{右端相遇}与\emph{左端相遇}
之间的时间间隔为 $\Delta t$。
求另一根杆相对于该杆的速率 $v$。
\exsolution
在第一杆静止系中，它长 $L_0$；
第二根杆沿轴运动且发生长度收缩，
长度为 $L=L_0\sqrt{1-\beta^2}$。
取两个端点的相遇顺序使移动杆在这段时间内
相对于第一杆扫过两杆长度之和，
因此
\[
v\Delta t=L_0+L
=L_0(1+\sqrt{1-\beta^2}).
\]
令 $a=L_0/(c\Delta t)$，
则 $\beta=a(1+\sqrt{1-\beta^2})$。
移项平方并剔除零速或不符合相遇顺序的分支，得
$\beta=2a/(1+a^2)$。
于是
\[
\boxed{v=\frac{2L_0/\Delta t}{
1+[L_0/(c\Delta t)]^2}}.
\]
\excheck 相遇方式必须是题定“右端相遇”和“左端相遇”
这两个不同事件；不能仅按一根杆的收缩长度
除以 $\Delta t$，会漏掉另一杆自身长度。
